The three bypass channels share a structural property: they do not require a majority of beneficiary legislators to vote against their interests. Each routes reform through a different institutional actor---voters, courts, or Congress---who does not face the same incumbency rent calculation.

\subsection{Channel 1: Direct Democracy}

Twenty-six states permit citizen ballot initiatives, which allow voters to enact law directly without legislative approval. This channel is the most frequently used bypass for redistricting reform. California's Proposition 11 (2008) created the Citizens Redistricting Commission, and Proposition 20 (2010) extended its authority to congressional districts---both over strong legislative opposition. Colorado's Amendment Y (congressional) and Amendment Z (state legislative) passed in 2018 with 71\% and 70\% support respectively. Michigan's Proposal 2 (2018) established an independent commission and passed with 61\% support.

The ballot initiative channel has three prerequisites. First, the state must permit constitutional amendments or statutory changes via initiative---20 states do not. Second, reformers must mount a successful petition drive (typically requiring signatures from 5--10\% of registered voters). Third, the measure must survive a statewide popular vote. The first two prerequisites filter out states with legislatively controlled initiative processes; the third relies on the gap between public opinion (supporting reform) and legislative behavior (opposing it).

The political economy of the initiative channel differs from the legislative channel in a crucial way: incumbents cannot vote against the initiative, but they can spend against it. California's redistricting reform battles saw millions of dollars in campaign spending from incumbent-aligned groups. Understanding the funding asymmetry---reformers typically outspent by incumbents---is essential for designing initiative strategies.

\subsection{Channel 2: Judicial Intervention}

Courts have invalidated redistricting maps on racial, partisan, and state-constitutional grounds. The judicial channel differs from the initiative channel in that it is reactive: courts act only when maps are challenged and only within their jurisdiction.

Federal courts currently offer limited relief for purely partisan claims. \textit{Rucho v. Common Cause}, 588 U.S.\ 684 (2019), held that federal courts lack a ``judicially manageable standard'' for evaluating partisan gerrymanders. However, state courts remain fully open. \textit{League of Women Voters v. Pennsylvania}, 178 A.3d 737 (Pa. 2018), struck the Pennsylvania congressional map under the state constitution. \textit{Harper v. Hall}, 380 N.C.\ 317 (2022), did the same in North Carolina. \textit{Harkenrider v. Hochul}, 38 N.Y.3d 494 (2022), invalidated New York's congressional map as an unconstitutional gerrymander.

Racial gerrymandering remains a federal judicial remedy. \textit{Alabama Legislative Black Caucus v. Alabama}, 575 U.S.\ 254 (2015), and \textit{Allen v. Milligan}, 599 U.S.\ 1 (2023), demonstrate that Section 2 of the Voting Rights Act continues to require congressional maps to preserve minority voting strength. Courts have appointed special masters to draw remedial maps in Alabama, North Carolina, New York, and Pennsylvania since 2010.

The judicial channel's principal limitation is speed: redistricting litigation routinely takes two to four years, meaning one or more election cycles pass under the challenged map before relief arrives. The \texttt{bisect} platform addresses this limitation directly: because algorithmic maps are byte-reproducible from public inputs, courts can order immediate remediation with a verifiable, neutral baseline. This point is developed in depth in P.4.

\subsection{Channel 3: Federal Legislation}

Congress possesses authority under Article I, \S4 to regulate the ``Manner'' of congressional elections, including the drawing of district lines. This power is not merely theoretical: Congress has previously mandated single-member districts (2 U.S.C.\ \S2c) and has regulated ballot access, campaign finance, and voting procedure. The precedent closest to redistricting is 2 U.S.C.\ \S2a, the Huntington-Hill apportionment statute, which mandates a specific mathematical algorithm for allocating seats among states.

Federal legislation is the most powerful bypass channel because it preempts state law and reaches all 50 states simultaneously. It is also the most difficult: it requires bicameral majorities, presidential signature, and sufficient political will to overcome filibuster or other blocking mechanisms. The Districting Integrity Act (DIA), analyzed in detail in P.1, is the model federal statute developed in conjunction with this research program.

The federal channel is further complicated by \textit{Rucho}'s suggestion that legislative rather than judicial solutions are appropriate for partisan gerrymandering---which is an implicit endorsement of the federal channel that reformers have not yet fully exploited. The case's language emphasizes that ``[f]ederal courts are not the only potential check on partisan gerrymandering.'' This rhetorical structure invites Congress to act.

\subsection{Interactions Among Channels}

The three channels are not independent. Judicial decisions can accelerate legislative reform by raising the cost of non-compliance. Initiative campaigns can create political pressure that makes legislative reform more likely. Federal legislation can preempt state-level debates, providing a national floor below which states cannot fall. The most successful reform ecosystems combine elements of all three: California's commission was created by initiative (Channel 1) but has been reinforced by judicial decisions interpreting the initiative's criteria (Channel 2) and would be backstopped by federal legislation (Channel 3) if the state legislature attempted to reassert control.
